\providecommand{\BM}{\textbf{[B]}}
\providecommand{\KM}{\textbf{[K]}}
\providecommand{\SM}{\textbf{[S]}}
\providecommand{\XM}{\textbf{[X]}}

\chapter*{The Yukawa Mechanism as a Consequence of $\tilde T\cdot m=1$}
\addcontentsline{toc}{chapter}{The Yukawa Mechanism as a Consequence of $\tilde T\cdot m=1$}

{\small\noindent\textbf{Abstract.}\quad\ignorespaces
In the Standard Model the Yukawa coupling $y_i = m_i/v$ is an empirical
assumption: it is chosen so that the observed masses emerge after
electroweak symmetry breaking. This document shows that the same coupling
follows from the FFGFT ground relation $\tilde T\cdot m = 1$ once the
electroweak scale $v$ is allowed to depend on position. The
mass-proportional coupling of all fermions to the Higgs is then not an
assumption but the fluctuating form of $\tilde T\cdot m = 1$. The Higgs
particle $h(x)$ is the vibrational quantum of the lattice scale --- not an
independent scalar field with an unknown potential. Its mass is the
residual stiffness of the lattice: $M_W^2+M_Z^2+m_h^2 = v^2/2$ (trace rule,
$0.45\,\%$). Together with $\sin^2\theta_W = 2/9$ and $m_h/M_Z = 11/8$,
$M_Z$, $M_W$ and $m_h$ follow from $v$ alone to within $0.2$--$0.3\,\%$.
}\medskip

\section*{Notation}
\addcontentsline{toc}{section}{Notation}

\medskip
\begin{tabular}{@{}ll@{}}
\textbf{[B]} & algebraically proven \\
\textbf{[K]} & numerically confirmed \\
\textbf{[S]} & structurally plausible, derivation open \\
\textbf{[X]} & tested and negative \\
\end{tabular}

\medskip
\begin{tabular}{@{}lp{9cm}@{}}
\toprule
\textbf{Symbol} & \textbf{Meaning} \\
\midrule
$\tilde T_i$ & proper-time scaling of mode $i$ \\
$m_i$ & rest mass of particle $i$ \\
$v$ & electroweak vacuum scale ($246.22$\,GeV) \\
$h(x)$ & physical Higgs field (fluctuation of $v$) \\
$y_i$ & Yukawa coupling (SM definition: $m_i = y_i v$) \\
$r_i$, $p_i$ & prefactor and exponent of the FFGFT mass formula \\
$\xi$ & FFGFT scaling parameter ($4/30000$) \\
\midrule
$\psi_i$, $\bar\psi_i$ & fermion field of particle $i$ and its Dirac adjoint \\
$\mathcal L$ & Lagrangian density \\
$H$ & complex lattice amplitude (Higgs field), $H = (v+h)/\sqrt2$ \\
$n_{\phi,i}$, $n_{\theta,i}$ & winding numbers on the $T^4$ torus \\
$n_c$ & colour multiplicity \\
$d_{{\rm akt},i}$ & number of active torus dimensions (Doc.~189) \\
\midrule
$M_W$, $M_Z$ & masses of the $W$ and $Z$ bosons (pole masses) \\
$m_h$, $m_t$ & mass of the Higgs particle, of the top quark \\
$m_s$ & mass of the strange quark \\
$W^\pm$ & charged gauge boson; one complex mode, $W^\pm = (W^1\mp iW^2)/\sqrt2$ \\
$W^3$, $B$ & neutral gauge fields before mixing ($SU(2)$, $U(1)$) \\
$\gamma$ & photon \\
$\theta_W$ & Weinberg angle, on-shell $\sin^2\theta_W = 1 - M_W^2/M_Z^2$ \\
$g$, $g'$ & gauge couplings of $SU(2)$ and $U(1)$ \\
$\lambda$ & Higgs self-coupling, $m_h^2 = 2\lambda v^2$ \\
$y_t$ & Yukawa coupling of the top quark \\
$G_F$ & Fermi constant \\
$\tau_\mu$ & muon lifetime (the only measured quantity behind $v$) \\
\midrule
$\phi$ & golden ratio $(1+\sqrt5)/2$ \\
$Q$ & Koide quotient $\sum m_i/(\sum\sqrt{m_i})^2$ \\
$\alpha$, $\alpha_s$ & fine-structure constant, strong coupling \\
$\Lambda_{\rm QCD}$ & confinement scale of the strong interaction \\
$R_{\rm Proton}$ & charge radius of the proton \\
$\theta_{12}$, $\theta_{13}$, $\theta_{23}$ & neutrino mixing angles \\
$\delta_{CP}$ & CP-violating phase \\
$K_{\rm frak}$ & fractal correction factor $1-100\xi$ (Doc.~133) \\
$\mathrm{GF}(q)$, $\mathrm{GF}(q)^*$ & Galois field with $q$ elements, its multiplicative group \\
$A_5$ & icosahedral group (60 elements) \\
\bottomrule
\end{tabular}

\section{The Ground Relation and Its SM Notation}

FFGFT sets for every fermion mode $i$:
\begin{equation}
  \tilde T_i \cdot m_i = 1, \qquad
  m_i = r_i\,\xi^{p_i}\,v.
  \label{eq:T0}
\end{equation}
The Yukawa coupling of the Standard Model is defined by $m_i = y_i\,v$,
hence:
\begin{equation}
  \boxed{y_i = r_i\,\xi^{p_i}.}
  \label{eq:yuk_ffgft}
\end{equation}
This is not a matching but a direct rewriting \BM. Here (Doc.~189):
\begin{align}
  r_i &= \frac{n_{\phi,i}^2}{n_{\theta,i}\,n_c},
  & p_i &= \frac{d_{{\rm akt},i}}{3},
\end{align}
where $n_{\phi,i}$, $n_{\theta,i}$ are the winding numbers on the $T^4$
torus and $n_c$ is the colour multiplicity. The Yukawa coupling is thus
fully determined by the torus topology --- it has no free parameter \BM.

\section{Derivation of the Yukawa Vertex from $\tilde T\cdot m = 1$}

\subsection{The vacuum case}

In the vacuum $v = {\rm const}$. Eq.~\eqref{eq:T0} gives
$\tilde T_i = 1/m_i$ --- every mode has a fixed proper-time scaling.
The Lagrangian mass term reads:
\begin{equation}
  \mathcal L \supset -m_i\,\bar\psi_i\psi_i.
\end{equation}

\subsection{Fluctuating lattice scale}

Now allow $v$ to fluctuate with position:
\begin{equation}
  v(x) = v + h(x), \qquad |h(x)| \ll v.
\end{equation}
The ground relation $\tilde T_i\cdot m_i = 1$ must continue to hold --- not
as a global normalisation but as a local bond between proper time and mass.
This means:
\begin{equation}
  m_i(x) = \frac{1}{\tilde T_i} = r_i\,\xi^{p_i}\,v(x)
  = r_i\,\xi^{p_i}\,\bigl(v + h(x)\bigr)
  = m_i^{(0)}\!\left(1 + \frac{h(x)}{v}\right).
  \label{eq:m_local}
\end{equation}

\subsection{The coupling term}

Inserting into the mass term gives \BM:
\begin{equation}
  \mathcal L \supset -m_i(x)\,\bar\psi_i\psi_i
  = -m_i^{(0)}\,\bar\psi_i\psi_i
    \;-\; \frac{m_i^{(0)}}{v}\,\bar\psi_i\psi_i\,h(x).
  \label{eq:yukawa_derived}
\end{equation}
The second term is exactly the Yukawa vertex of the Standard Model with
$y_i = m_i^{(0)}/v$. Here it does not follow from an assumption about the
coupling structure but from the single FFGFT relation $\tilde T\cdot m = 1$
under spatial fluctuations of $v$.

\section{Mass-Proportional Coupling as a Consequence}

Eq.~\eqref{eq:yukawa_derived} holds for \emph{all} fermions $i$
\emph{simultaneously}, with the same $v$ and the same $h(x)$. This means:

\begin{quote}
\textbf{[B]} The mass-proportional coupling $y_i = m_i/v$ of all fermions to
the Higgs boson is a necessary consequence of $\tilde T\cdot m = 1$ under
spatial fluctuation of the lattice scale. It is not a free assumption of the
model.
\end{quote}

This is the content of the Yukawa mechanism in FFGFT language: not
``particles couple to a scalar field'' but ``mass is a local property of the
lattice, and $h(x)$ is the oscillation of this property''.

\section{Interpretation of the Higgs Particle}
\label{sec:interp}

In this derivation the field $h(x)$ is the position-dependent deviation of
the lattice scale $v$ from its vacuum value. The physical Higgs particle ---
the resonance at $125.2$\,GeV --- is then the vibrational quantum of this
deviation.

This differs conceptually from the Standard Model:

\begin{center}
\small\begin{tabular}{@{}lll@{}}
\toprule
 & \textbf{Standard Model} & \textbf{FFGFT} \\
\midrule
Origin of mass & Yukawa coupling (postulated) & $\tilde T\cdot m = 1$ (ground relation) \\
Higgs field & independent scalar field & fluctuation of $v(x)$ \\
Coupling $y_i = m_i/v$ & empirical parameter & derived from Eq.~\eqref{eq:m_local} \\
Number of free parameters & $n_{\rm fermions}$ & $0$ \\
Higgs potential $V(H)$ & Mexican hat (postulated) & residual stiffness of the lattice (§\ref{sec:spur}) \\
Higgs mass $m_h$ & free parameter & trace rule (§\ref{sec:spur}) \KM \\
\bottomrule
\end{tabular}
\end{center}

The derivation so far shows that $h(x)$ \emph{must} exist and \emph{must}
couple with $m_i/v$. The restoring force of the oscillation --- the
potential $V(h)$ and hence the Higgs mass --- is treated in
§\ref{sec:higgsmasse} and §\ref{sec:spur}.

\section{Higgs Mass from the Galois Chain \KM}
\label{sec:higgsmasse}

Potential, Higgs mass and oscillation frequency are a single problem. This
section shows that a numerically sharp candidate exists that is built
entirely from Galois structures of the corpus.

\paragraph{The 11/8 chain.}
The three quantities $M_Z$, $m_h$, $m_t$ satisfy:
\begin{equation}
  M_Z : m_h : m_t \;=\; 1 : \tfrac{11}{8} : \bigl(\tfrac{11}{8}\bigr)^2,
  \qquad
  m_h = M_Z\cdot\tfrac{11}{8}, \quad m_t = M_Z\cdot\bigl(\tfrac{11}{8}\bigr)^2.
  \label{eq:kette}
\end{equation}
Numerically \KM:
\begin{center}
\small\begin{tabular}{@{}lrrr@{}}
\toprule
Quantity & measured & FFGFT (Eq.~\eqref{eq:kette}) & deviation \\
\midrule
$M_W$ & $80.369$\,GeV & $M_Z\sqrt{7/9} = 80.420$\,GeV & $0.06\,\%$ \\
$M_Z$ & $91.188$\,GeV & input & --- \\
$m_h$ & $125.200$\,GeV & $M_Z\cdot11/8 = 125.383$\,GeV & $0.15\,\%$ \\
$m_t$ & $172.570$\,GeV & $M_Z\cdot(11/8)^2 = 172.402$\,GeV & $0.10\,\%$ \\
\bottomrule
\end{tabular}
\end{center}
All four heavy bosons and fermions follow from two Galois numbers plus $v$.

\paragraph{Galois origin of the numbers.}
Both factors are field quantities of the corpus:
$8 = |\mathrm{GF}(9)^*| = 3^2 - 1$ (Doc.~336) and
$11 \in \text{orbit}\{7,11,21\}$ under $x\mapsto 3x \bmod 26$
in $\mathrm{GF}(27)^*$ (Doc.~340, the same structure as $\Delta m^2_{\rm atm}$).
The Weinberg factor $\sqrt{7/9}$ comes from the $A_5$ matrix element $2/9$
(Doc.~372). The chain thus links three different Galois structures into
one consistent numerical picture \KM.

\paragraph{Classification.}
Eq.~\eqref{eq:kette} is a numerically confirmed relation; the Galois origin
of the numbers 11 and 8 is a reading \SM. A derivation of the Higgs mass
that does without these numbers is provided by the trace rule in
§\ref{sec:spur}; together the two fix the entire electroweak boson spectrum
from $v$ alone.

\section{The Trace Rule: Higgs Mass as Residual Stiffness of the Lattice}
\label{sec:spur}

\paragraph{Picture.}
Coupled resonant circuits detune each other: the coupling pushes the
eigenfrequencies apart. What is conserved in the process is the sum of the
squared frequencies --- the trace of the stiffness matrix. In this picture
the Higgs mass is not a resonance of its own with its own potential but the
share of the lattice stiffness left over by the gauge bosons.

\paragraph{The rule.}
\begin{equation}
  M_W^2 + M_Z^2 + m_h^2 \;=\; |\langle H\rangle|^2 \;=\; \frac{v^2}{2},
  \qquad
  m_h^2 = \frac{v^2}{2} - M_W^2 - M_Z^2 .
  \label{eq:spur}
\end{equation}

\paragraph{The four building blocks.}
\begin{enumerate}
\item \textbf{Conservation under coupling \BM.} The sum of the eigenvalues
  of a symmetric matrix is its trace, independent of the basis. The
  electroweak mixing is a pure rotation by $\theta_W$: before mixing, $W^3$
  and $B$ carry the stiffnesses $g^2v^2/4$ and $g'^2v^2/4$, afterwards the
  photon ($0$) and the $Z$ ($M_Z^2$) --- the sum is identical, the mixing
  matrix orthogonal.
\item \textbf{No intrinsic stiffness of the coupling \BM.} A mechanical
  coupling spring is a component of its own and adds stiffness to the
  diagonal. The vacuum is empty; the coupling can only redistribute. That is
  exactly the property of a rotation (point~1).
\item \textbf{Counting once per mode.} $W^+$ and $W^-$ are one complex mode
  with two senses of circulation: in the corpus the antiparticle is the same
  field with reversed sign (Doc.~042, 053), and the trace runs over $m^2$,
  where the sign drops out. Consistently, Koide counts three lepton masses,
  not six.
\item \textbf{Budget $= v^2/2$ \BM.} The Higgs \emph{field} is the ground
  lattice itself; the total stiffness belongs to it and is the squared
  modulus of its vacuum amplitude. Since the amplitude is complex (modulus
  and phase), the canonical normalisation of the real oscillation $h$
  requires the factor $1/\sqrt2$: $H = (v+h)/\sqrt2$, hence
  $|\langle H\rangle|^2 = v^2/2$. It is the same $\sqrt2$ as in
  $W^\pm = (W^1\mp iW^2)/\sqrt2$ --- the argument that counts $W$ once also
  fixes the factor $1/2$.
\end{enumerate}
The Higgs \emph{particle} $h$ is the eigen-oscillation of the lattice; its
stiffness $m_h^2$ is the remainder in Eq.~\eqref{eq:spur}. The
Mexican-hat potential of the Standard Model is therefore not postulated:
its curvature at the minimum is this residual stiffness.

\paragraph{Numbers \KM.}
With $M_W^2 = \tfrac79 M_Z^2$ (Doc.~372) it follows that
$m_h = \sqrt{v^2/2 - \tfrac{16}{9}M_Z^2} = 124.6$\,GeV
(measured $125.2$\,GeV, $-0.47\,\%$). The balance itself holds to
$0.45\,\%$.

\paragraph{Boson spectrum from $v$ alone \KM.}
Combining Eq.~\eqref{eq:spur} with $\sin^2\theta_W = 2/9$ and the chain
$m_h = \tfrac{11}{8}M_Z$ (§\ref{sec:higgsmasse}) leaves
\[
  M_Z^2\Bigl(\tfrac79 + 1 + \tfrac{121}{64}\Bigr) = \frac{v^2}{2}
  \quad\Longrightarrow\quad
  \frac{M_Z^2}{v^2} = \frac{288}{2113},
\]
and all three boson masses follow from $v$ without any measured mass as
input:
\begin{center}
\small\begin{tabular}{@{}lrrr@{}}
\toprule
 & FFGFT & measured & deviation \\
\midrule
$M_Z$ & $90.90$\,GeV & $91.19$\,GeV & $-0.31\,\%$ \\
$M_W$ & $80.17$\,GeV & $80.37$\,GeV & $-0.25\,\%$ \\
$m_h$ & $124.99$\,GeV & $125.20$\,GeV & $-0.17\,\%$ \\
\bottomrule
\end{tabular}
\end{center}
As couplings: $g = 0.6512$ (on-shell $2M_W/v = 0.6528$),
$\lambda = 0.1288$ (from $m_h$: $0.1293$); the Higgs self-coupling is fixed
by $4g^2/7 + 2\lambda = 1/2$ and is no longer a free parameter.

\paragraph{Accuracy.}
All three values lie $0.2$--$0.3\,\%$ low, in the same direction.
Eq.~\eqref{eq:spur} is a tree-level relation; the electroweak radiative
corrections are of the order of a few per cent. Without a loop calculation
no better agreement than the one reached here can be expected. The choice
of $v$ shifts the result: with the SM definition $246.22$\,GeV the balance
holds to $+0.45\,\%$, with the FFGFT values from Doc.~006/R104 ($248.9$ and
$245.6$\,GeV) to $-1.7$ and $+1.0\,\%$ respectively. Which $v$ is the right
one is the question registered in R104/R112 (convention-free only via
$\tau_\mu$), not a gap in the trace rule.

\paragraph{Fermion side: the trace rule applies to lattice modes only.}
The obvious extension ``the top carries the same budget'',
$m_t^2 = v^2/2$ ($y_t = 1$), is incompatible with the chain: together with
Eq.~\eqref{eq:spur} it requires $m_t/M_Z = \sqrt{2113/576} = 1.915$, while
the measured value is $1.8925$ ($+1.2\,\%$) and $(11/8)^2 = 1.8906$ agrees to
$0.10\,\%$. The reading $y_t = 1$ is discarded \XM: it has no justification
beyond the analogy, and it is the worse one. In addition, the top mass
itself is scheme-dependent --- pole mass $172.6$\,GeV, $\overline{\rm MS}$
mass $162.5$\,GeV, a spread of $5.8\,\%$, corresponding to $y_t = 0.99$ and
$0.93$. Quarks are not free asymptotic modes (Doc.~372 §4); the trace rule
applies to the eigen-oscillations of the lattice ($W$, $Z$, $h$), not to
confined fermions. For the top mass the chain $m_t = (11/8)^2M_Z$ applies
\KM.

\paragraph{Rational ratios versus $v$-dependent relations.}
The accuracy of the relations is ordered by whether they contain $v$:
\begin{center}
\small\begin{tabular}{@{}llr@{}}
\toprule
Relation & contains & deviation \\
\midrule
$M_W^2/M_Z^2 = 7/9$ & pole masses only & $0.06\,\%$ \\
$m_t/M_Z = (11/8)^2$ & pole masses only & $0.10\,\%$ \\
$m_h/M_Z = 11/8$ & pole masses only & $0.15\,\%$ \\
$M_W^2+M_Z^2+m_h^2 = v^2/2$ & $v$ & $0.45\,\%$ \\
$m_t^2 = v^2/2$ (discarded) & $v$ & $1.2\,\%$ \\
\bottomrule
\end{tabular}
\end{center}
Ratios of measured pole masses are convention-free; this is where the
rational Galois numbers agree best --- as is to be expected, since they are
exact numbers, not parameters. $v$, by contrast, is not a measurement but
the SM definition $(\sqrt2\,G_F)^{-1/2}$, whose irrational factors stem from
the SM formulation and which carries a theory correction of $-0.44\,\%$ from
muon decay (Doc.~351, R112). Every relation containing $v$ inherits these
conventions. The residual of the trace rule ($0.45\,\%$) is of the size of
this correction; that it stems from the definition of $v$ and not from the
rule is a testable conjecture \SM.

\section{Inventory of the Particle Zoo (as of 21~September 2026)}
\label{sec:bestand}

The following overview summarises what the corpus, after Doc.~372 and this
document, delivers for each sector. In each case the best derivation is
decisive; older, superseded formulas are recorded in the register
(Doc.~190).

\begin{center}
\small
\begin{tabular}{@{}>{\raggedright\arraybackslash}p{2.2cm}>{\raggedright\arraybackslash}p{6.6cm}>{\raggedright\arraybackslash}p{2.3cm}l@{}}
\toprule
\textbf{Sector} & \textbf{Quantity} & \textbf{Source} & \textbf{Status} \\
\midrule
Charged & mass ratios $e:\mu:\tau$ from the matrix element $2/9$ & 292, 293, 367 & \BM/\KM \\
leptons & $\phi$ skeleton of the weights, $4/3 = 2Q$ & 368, 371 & \BM \\
\midrule
Neutrinos & masses from the $\mathrm{GF}(27)^*$ orbits & 340 & \KM \\
 & $\nu_3$ Dirac, $\nu_{1,2}$ Majorana & 343 & \BM \\
 & $\sin^2\theta_{13} = (25\xi)^{2/3}$ & 340 & \KM \\
 & $\sin^2\theta_{23}\in\{4/9,5/9\}$, $|\sin^2\theta_{23}-\tfrac12| = \tfrac1{18}$ & 372 & \KM \\
 & $\theta_{12}$ (order of magnitude only), $\delta_{CP}$ & 340 & \SM \\
 & absolute scale $m_\nu$ (ansatz) & 007 & \SM \\
\midrule
Gauge & photon and eight gluons massless & 320, 321, 339 & \BM \\
bosons & $M_W^2/M_Z^2 = 7/9$ ($\sin^2\theta_W = 2/9$ on-shell) & 372 & \KM \\
 & $M_Z^2/v^2 = 288/2113$ ($M_Z$ from $v$ alone) & 373 & \KM \\
\midrule
Higgs & Yukawa coupling $y_i = m_i/v$ from $\tilde T\cdot m = 1$ & 373 & \BM \\
 & $m_h = \tfrac{11}{8}M_Z$ & 373 & \KM \\
 & trace rule $M_W^2+M_Z^2+m_h^2 = v^2/2$, self-coupling fixed & 373 & \KM \\
\midrule
Top & $m_t = (\tfrac{11}{8})^2 M_Z$ & 373 & \KM \\
\midrule
Other & exponents $p_i = d_{\rm akt}/3$ & 189 & \KM \\
quarks & prefactors $r_i$ in the Galois grid; derivation open & 006, 358 & \BM/\SM \\
 & unrolling as for leptons ($Q = 2/3$) & 372 & \XM \\
\midrule
Hadrons & confinement as a topological boundary condition & 319 & \SM \\
 & Gell-Mann--Okubo splitting consistent with $m_s$ & 372 & \KM \\
 & $R_{\rm Proton}$ structurally classified & 372 & \SM \\
 & $\Lambda_{\rm QCD}$ from the $\xi$ ladder & 372 & \XM \\
\midrule
Couplings & $\alpha$ from Galois orders & 338 & \KM \\
 & $\alpha_s$ via the Galois running & 160, R78 & \KM \\
 & Cabibbo angle & 372 & \SM \\
\midrule
Scale & $v$ from the lattice & 006, R104 & \KM \\
 & convention-free test only via $\tau_\mu$ & 351, R112 & \SM \\
\bottomrule
\end{tabular}
\end{center}

\paragraph{Ordering principle.}
Precise are the free modes --- the charged leptons --- and the oscillations
of the lattice itself ($W$, $Z$, $h$). Blurred is everything confined:
quarks and hadrons, whose masses depend on the scheme. And the exact Galois
ratios between pole masses agree better than any relation that contains the
SM quantity $v$ (§\ref{sec:spur}).

\paragraph{Open.}
$\theta_{12}$ and $\delta_{CP}$; the quark prefactors $r_i$ apart from the
top; $\Lambda_{\rm QCD}$; the Cabibbo angle; the geometric derivation of the
hadronic factor in Doc.~160 (R123); the convention-free determination of $v$
via the muon lifetime (R112).

\section{Result}

\begin{enumerate}
\item The Yukawa coupling $y_i = r_i\,\xi^{p_i}$ is a direct rewriting of
  the torus topology --- not a free assumption \BM.
\item The Yukawa vertex $-(m_i/v)\,\bar\psi_i h\psi_i$ follows from
  $\tilde T\cdot m = 1$ under $v\to v+h(x)$ for all fermions
  simultaneously \BM.
\item The mass-proportional coupling (the fingerprint of the Higgs
  discovery) is a consequence, not a parameter \BM.
\item $h(x)$ is the vibrational quantum of the lattice scale. The Higgs
  potential is not a separate postulate: its curvature at the minimum is the
  residual stiffness of the lattice (§\ref{sec:spur}).
\item \textbf{Higgs mass candidate} \KM: $m_h = M_Z\cdot11/8 = 125.383$\,GeV
  to $0.15\,\%$; $m_t = M_Z\cdot(11/8)^2$ to $0.10\,\%$. Both factors ---
  $11\in\mathrm{GF}(27)^*$, $8=|\mathrm{GF}(9)^*|$ --- are Galois quantities
  of the corpus; their origin is a reading \SM.
\item \textbf{Trace rule}: $M_W^2+M_Z^2+m_h^2 = v^2/2$ to $0.45\,\%$ \KM.
  Conservation under coupling, absence of intrinsic stiffness, counting and
  the factor $1/2$ are justified \BM; $m_h = \sqrt{v^2/2-\tfrac{16}{9}M_Z^2}$.
\item \textbf{Boson spectrum from $v$} \KM: $M_Z^2/v^2 = 288/2113$;
  $M_Z$, $M_W$, $m_h$ to $0.17$--$0.31\,\%$, at tree level.
\item Inventory of the entire particle zoo in §\ref{sec:bestand}.
\item Top: the trace rule applies to lattice modes only; $y_t = 1$ discarded
  \XM, $m_t = (11/8)^2M_Z$ \KM.
\item Accuracy is ordered by convention dependence: rational pole-mass
  ratios $0.06$--$0.15\,\%$, relations with $v$ $0.45\,\%$. Open \SM: whether
  the residual of the trace rule stems from the definition of $v$
  (R104/R112).
\end{enumerate}

\medskip
\noindent\textbf{Check script:} \texttt{2/python/Dok373\_Skripte/pruef\_373\_yukawa.py}

\medskip
\noindent\textbf{Register (Doc.~190):} R122 updated (Higgs mass: candidates available).

\begin{thebibliography}{9}
\bibitem{dok006} J.~Pascher, \emph{Doc.~006 --- T0 particle masses} (2025).
\bibitem{dok189} J.~Pascher, \emph{Doc.~189 --- Torus derivation of the winding numbers} (2026).
\bibitem{dok372} J.~Pascher, \emph{Doc.~372 --- The matrix-element principle} (2026).
\end{thebibliography}
